%%
%% winter-lecture27.tex
%% 
%% Made by Alex Nelson <pqnelson@gmail.com>
%% Login   <alex@lisp>
%% 
%% Started on  2026-03-07T15:21:18-0800
%% Last update 2026-03-07T15:21:18-0800
%% 

\lecture{}

\begin{definition}
A topological space $Z$ is a \define{Baire space} if every countable
intersection of open dense subsets of $Z$ is dense in
$Z$. Equivalentls, every countable union of nowhere dense subsets has
empty interior.
\end{definition}

\begin{fact}
Every Banach space is a Baire space.
\end{fact}

\begin{lemma}[Approximation lemma]
If $B_{Y}(0,r)\subset\closure{T(B_{X}(0,1))}$, then
$B_{Y}(0,r/2)\subset T(B_{X}(0,1))$.
\end{lemma}

\begin{remark}
Completeness of $X$ is the key point. We will build correction terms
in $X$ where norms form a geometric series $\sum2^{-j}$, their sum
stays in $X$ and error terms go to zero.
\end{remark}

\begin{proof}[Proof (Approximation lemma)]
Fix $y\in B_{Y}(0,r/2)$. We will construct a series
\begin{equation}
\sum^{\infty}_{n=0}x_{n}=x
\end{equation}
with $\|x\|<1$ and $Tx=y$. For partial sums we write
\begin{equation}
s_{n}=\sum^{n}_{k=0}x_{k},
\end{equation}
with $e_{n}=y-Ts_{n}$ is the error after $n+1$ corrections.

\textsc{Step 1}: By scaling the hypothesis, $B_{Y}(0,r)\subset\closure{T(B_{X}(0,1))}$,
which implies $B_{Y}(0,r/2)\subset\closure{T(B_{X}(0,1/2))}$. Since
$y\in B_{Y}(0,r/2)$, we choose $x_{0}\in B_{X}(0,1/2)$ such that
\begin{equation}
\|y-Tx_{0}\|<\frac{r}{4}.
\end{equation}
This says $\|e_{0}\|<r/4$.

\textsc{Inductive Step}: Suppose we have $x_{0}$, \dots, $x_{n}$ which
have been chosen such that $\|e_{n}\|<r/2^{n+2}$. Scaling the hypothesis
$B_{Y}(0,r/2^{n+2})\subset\closure{T(B_{X}(0,1/2^{n+2}))}$. Since
$e_{n}$ lies in this ball, we have chosen $e_{n}\in B_{X}(0,2^{-(n+2)})$
with
\begin{equation}
\|e_{n}-Tx_{n+1}\|<r/2^{n+3}.
\end{equation}
If
\begin{equation}
s_{n+1}:=s_{n}+x_{n+1},
\end{equation}
then
\begin{equation}
e_{n+1}=y_{n}-Ts_{n+1}=e_{n}-Tx_{n+1},
\end{equation}
so $\|e_{n+1}\|<r/2^{n+3}$. From the construction,
\begin{equation}
\|x_{n}\|<2^{-(n+1)}\quad\mbox{for }n\geq0.
\end{equation}
Hence
\begin{equation}
\sum^{\infty}_{n=0}\|x_{n}\|<\sum^{\infty}_{n=0}2^{-n-1}=1.
\end{equation}
Since $X$ is Banach, this series $\sum x_{n}$ converges to some $x\in
X$. Also,
\begin{equation}
\|x\|\leq\sum^{\infty}_{n=0}\|x_{n}\|<1,
\end{equation}
so $x\in B_{X}(0,r=1)$.

Now, $s_{n}\to x$ in the $X$ norm, so continuity gives us $Ts_{n}\to Tx$.
On the other hand,
\begin{subequations}
  \begin{align}
\|y-Ts_{n}\| &= \|e_{n}\|\\
&<\frac{r}{2^{n+2}}\xrightarrow{n\to\infty}0,
  \end{align}
\end{subequations}
so $Ts_{n}\to y$ and therefore
\begin{equation}
Tx=y.
\end{equation}
We have shown $y\in T(B_{X}(0,r=1))$. Since $y\in B_{Y}(0,r/2)$
was arbitrary, $B_{Y}(0,r/2)\subset T(B_{X}(0,r=1))$ as desired.
\end{proof}

\begin{theorem}[Open mapping theorem]
Let $X$ and $Y$ be Banach spaces. Let $T\colon X\to Y$ be a surjective bounded
linear operator. Then $T$ is an open map.
\end{theorem}

\begin{proof}
By the Approximation Lemma, there exists a $\delta>0$ such that
$B_{Y}(0,\delta)\subset T(B_{X}(0,1))$. Let $U\subset X$ be an
arbitrary open set, let $y\in T(U)$, then consider $x\in U$ such that $Tx=y$
[such an $x$ exists since $T$ is surjective]. Since $U$ is open, there
exists $\varepsilon>0$ such that $B_{X}(x,\varepsilon)\subset U$. By linearity,
\begin{subequations}
  \begin{align}
T(B_{X}(x,\varepsilon))
&=T(x+B_{X}(0,\varepsilon))\\
&=Tx + T(B(0,\varepsilon))\\
&=Tx + \varepsilon T(B(0,1))\\
&=y + \varepsilon T(B(0,1)).
  \end{align}
\end{subequations}
Then $T(B_{X}(x,\varepsilon))\supset y+\varepsilon B_{Y}(0,\delta)$
since $B_{Y}(0,\delta)\subset T(B_{X}(0,1))$; i.e.,
\begin{equation}
B_{Y}(y,\varepsilon\delta)\subset T(B_{X}(x,\varepsilon)).
\end{equation}
Since $B_{X}(x,\varepsilon)\subset U$ we obtain
\begin{equation}
B_{Y}(y,\varepsilon\delta)\subset T(U).
\end{equation}
Since $y\in T(U)$ was arbitrary, we do this for every such $y\in T(U)$,
which gives us an open ball $B_{Y}(y,\widetilde{\varepsilon})\subset T(U)$
so $T(U)$ must be open.
\end{proof}

\begin{theorem}[Bounded inverse theorem]
Let $X$ and $Y$ be Banach spaces. If $T\colon X\to Y$ is a bijective
bounded linear operator, then $T^{-1}\colon Y\to X$ is automatically bounded.
\end{theorem}

\begin{proof}
Since $T$ is a surjective bounded linear operator between Banach
spaces, by the open mapping theorem we see $T$ is open. Let $U\subset X$
be an arbitrary open subset. Then
\begin{equation}
(T^{-1})^{-1}(U)=T(U)
\end{equation}
is open in $Y$ (since $T$ is an open map), so $T^{-1}$ is continuous
(since the preimage of open sets in $Y$ are open in $X$ under
$T^{-1}$). Hence the claim.
\end{proof}

\begin{example}
Let $X=c_{c}(\NN)$ be the set of compactly supported sequences:
$(a_{n})\in X$ iff $\exists N\in\NN\forall n\in\NN\ldotp n\geq N\implies a_{n}=0$.
We equip $X$ with the norm $\|(a_{n})\|=\sup_{n}|a_{n}|$.
Then $X$ is not complete. In particular, the sequence $a_{n}=1/n$ is
not in $X$ but it is the limit of the sequence
\begin{equation}
a^{(N)}_{n}=\begin{cases}
1/n & \mbox{if }n<N\\
0 & \mbox{if }n\geq N.
\end{cases}
\end{equation}
Then
\begin{equation}
\lim_{N\to\infty}(a^{(N)}_{n})=(a_{n}).
\end{equation}
Consider $T\colon X\to X$ defined by $T(a_{n})=(\frac{1}{n}a_{n})$.
This is a linear operator. We see it is bounded because
$\|Ta\|_{\infty}\leq\|a\|_{\infty}$ so the operator norm $\|T\|_{\text{op}}\leq1$.
We see $T$ is bijective on $c_{c}(\NN)$ and we find
\begin{equation}
T^{-1}((b_{1},b_{2},b_{3},\dots))=(b_{1},2b_{2},3b_{3},\dots,nb_{n},\dots).
\end{equation}
Then $\|T^{-1}\|_{\text{op}}=\infty$. So the bounded inverse theorem
\emph{fails} and we only removed the ``Banach'' property.
\end{example}